% Part: first-order-logic
% Chapter: sequent-calculus
% Section: proof-theoretic-notions

\documentclass[../../../include/open-logic-section]{subfiles}

\begin{document}

\begin{editorial}
  Bagean iki nglumpukake definisi relasi provabilitas lan konsistensi
  kanggo kalkulus sekuen. Cathetan koreksi sumber OLPL-005: editorial
  Inggris nyebut deduksi natural, nanging identitas berkas, notasi,
  lan kabeh derivasi ing bagian iki mratelakake kalkulus sekuen.
\end{editorial}

\iftag{FOL}
      {\olfileid{fol}{seq}{ptn}}
      {\olfileid{pl}{seq}{ptn}}

\olsection{Gagasan Teoretis-Bukti}

\begin{explain}
Kaya nalika nemtokake sawetara gagasan semantis wigati (validitas,
akibat, lan bisa dicukupi), saiki kita nemtokake \emph{gagasan
teoretis-bukti} kang cocog. Gagasan iki ora ditemtokake kanthi
nggunakake panyukupan !!{sentence}s ing !!{structure}s, nanging kanthi
nggunakake !!{derivability} utawa !!{nonderivability} sekuen-sekuen
tartamtu. Panemuan yen gagasan-gagasan kasebut padha iku panemuan
wigati. Kasunyatan manawa gagasan kasebut padha dadi isi
\emph{teorema kasahihan} lan \emph{teorema kelengkapan}.
\end{explain}

\begin{defn}[Teorema]
!!^a{sentence}~$!A$ iku \emph{teorema} yen ana !!a{derivation}
ing~$\Log{LK}$ saka sekuen $\quad \Sequent !A$. Kita nulis $\Proves
!A$ yen $!A$ iku teorema, lan $\Proves/ !A$ yen dudu teorema.
\end{defn}

\begin{defn}[!!^{derivability}]
!!^a{sentence}~$!A$ \emph{bisa !!{derivable} saka} himpunan
!!{sentence}s~$\Gamma$, ditulis $\Gamma \Proves !A$, yen lan mung yen
ana himpunan bagean winates~$\Gamma_0 \subseteq \Gamma$ lan urutan
$\Gamma_0'$ saka !!{sentence}s ing~$\Gamma_0$ saengga $\Log{LK}$
!!{derive}s $\Gamma_0' \Sequent !A$. Yen $!A$ ora bisa !!{derivable}
saka $\Gamma$, kita nulis $\Gamma \Proves/ !A$.
\end{defn}

Amarga aturan kontraksi, pelemahan, lan pertukaran, urutan lan cacah
!!{sentence}s ing~$\Gamma_0'$ ora wigati: yen sekuen $\Gamma_0'
\Sequent !A$ bisa !!{derivable}, mula semono uga $\Gamma_0'' \Sequent
!A$ kanggo saben $\Gamma_0''$ kang ngemot
!!{sentence}s kang padha karo~$\Gamma_0'$. Tuladhane, yen $\Gamma_0 =
\{!B, !C\}$, mula $\Gamma_0' = \tuple{!B, !B, !C}$ lan
$\Gamma_0'' = \tuple{!C, !C, !B}$ kalorone iku urutan kang mung ngemot
!!{sentence}s ing~$\Gamma_0$. Yen sekuen kang ngemot salah sijine bisa
!!{derivable}, sekuen kang ngemot sijine uga bisa, upamane:
\begin{prooftree}
  \AxiomC{}
  \Deduce$!B, !B, !C \fCenter !A$
  \RightLabel{\LeftR{\Contraction}}
  \UnaryInf$!B, !C \fCenter !A$
  \RightLabel{\LeftR{\Exchange}}
  \UnaryInf$!C, !B \fCenter !A$
  \RightLabel{\LeftR{\Weakening}}
  \UnaryInf$!C, !C, !B \fCenter !A$
\end{prooftree}
Wiwit saiki, yen $\Gamma_0$ iku himpunan winates saka !!{sentence}s,
kita bakal ngarani $\Gamma_0 \Sequent !A$ sembarang sekuen kang
antesedene arupa urutan !!{sentence}s ing~$\Gamma_0$, lan meneng-meneng
nglebokake kontraksi, pertukaran, lan pelemahan yen perlu.

\begin{defn}[Konsistensi]
Himpunan ukara~$\Gamma$ iku \emph{ora konsisten} yen lan mung yen ana
himpunan bagean winates~$\Gamma_0 \subseteq \Gamma$ saengga
$\Log{LK}$ !!{derive}s $\Gamma_0 \Sequent \quad$. Yen $\Gamma$ ora
kalebu ora konsisten, yaiku yen kanggo saben $\Gamma_0 \subseteq
\Gamma$ kang winates, $\Log{LK}$ ora !!{derive} $\Gamma_0 \Sequent
\quad$, kita ngarani himpunan iku \emph{konsisten}.
\end{defn}

\begin{prop}[Refleksivitas]
\ollabel{prop:reflexivity}
Yen $!A \in \Gamma$, mula $\Gamma \Proves !A$.
\end{prop}

\begin{proof}
Sekuen wiwitan $!A \Sequent !A$ bisa !!{derivable}, lan $\{!A\}
\subseteq \Gamma$.
\end{proof}
  
\begin{prop}[Monotonisitas]
\ollabel{prop:monotonicity}
Yen $\Gamma \subseteq \Delta$ lan $\Gamma \Proves !A$, mula $\Delta
\Proves !A$.
\end{prop}

\begin{proof}
Upamana $\Gamma \Proves !A$, yaiku ana $\Gamma_0 \subseteq \Gamma$
kang winates saengga $\Gamma_0 \Sequent !A$ bisa !!{derivable}.
Amarga $\Gamma \subseteq \Delta$, $\Gamma_0$ uga dadi himpunan bagean
winates saka~$\Delta$. Dadi, !!{derivation} saka $\Gamma_0 \Sequent
!A$ uga nuduhake $\Delta \Proves !A$.
\end{proof}

\begin{prop}[Transitivitas]
\ollabel{prop:transitivity}
Yen $\Gamma \Proves !A$ lan $\{!A\} \cup \Delta \Proves
!B$, mula $\Gamma \cup \Delta \Proves !B$.
\end{prop}

\begin{proof}
Yen $\Gamma \Proves !A$, ana $\Gamma_0 \subseteq \Gamma$ kang winates
lan !!a{derivation}~$\pi_0$ saka $\Gamma_0 \Sequent !A$. Yen
$\{!A\} \cup \Delta \Proves !B$, mula kanggo sawijining himpunan
bagean winates $\Delta_0 \subseteq \Delta$, ana
!!a{derivation}~$\pi_1$ saka $!A, \Delta_0 \Sequent !B$. Gatekna
!!{derivation} iki:
\begin{prooftree}
  \AxiomC{}
  \RightLabel{$\pi_0$}
  \Deduce$\Gamma_0 \fCenter !A$
  \AxiomC{}
  \RightLabel{$\pi_1$}
  \Deduce$!A, \Delta_0 \fCenter !B$
  \RightLabel{\Cut}
  \BinaryInf$\Gamma_0, \Delta_0 \fCenter !B$
\end{prooftree}
Amarga $\Gamma_0 \cup \Delta_0 \subseteq \Gamma \cup \Delta$, iki
nuduhake $\Gamma \cup \Delta \Proves !B$.
\end{proof}

Gatekna yen iki mligi ateges: yen $\Gamma \Proves !A$ lan $!A
\Proves !B$, mula $\Gamma \Proves !B$. Uga ndherek yen $!A_1,
\dots, !A_n \Proves !B$ lan $\Gamma \Proves !A_i$ kanggo saben~$i$,
mula $\Gamma \Proves !B$.

\begin{prop}
\ollabel{prop:incons}
$\Gamma$ ora konsisten yen lan mung yen $\Gamma \Proves {!A}$ kanggo
saben ukara~$!A$.
\end{prop}

\begin{proof}
Gladhen.
\end{proof}

\begin{prob}
Buktekna \olref[fol][seq][ptn]{prop:incons}
\end{prob}

\begin{prop}[Kekompakan]
  \ollabel{prop:proves-compact}
  \begin{enumerate}
  \item Yen $\Gamma \Proves !A$, mula ana himpunan bagean winates
    $\Gamma_0 \subseteq \Gamma$ saengga $\Gamma_0 \Proves !A$.
  \item Yen saben himpunan bagean winates saka~$\Gamma$ konsisten,
    mula $\Gamma$ konsisten.
  \end{enumerate}
\end{prop}

\begin{proof}
  \begin{enumerate}
    \item Yen $\Gamma \Proves !A$, mula ana himpunan bagean winates
      $\Gamma_0 \subseteq \Gamma$ saengga sekuen $\Gamma_0
      \Sequent !A$ nduweni !!a{derivation}. Akibate, $\Gamma_0
      \Proves !A$.
    \item Yen $\Gamma$ ora konsisten, ana himpunan bagean
      winates~$\Gamma_0 \subseteq \Gamma$ saengga $\Log{LK}$
      !!{derive}s $\Gamma_0 \Sequent \quad$. Nanging banjur $\Gamma_0$
      iku himpunan bagean winates saka~$\Gamma$ kang ora konsisten.
  \end{enumerate}
\end{proof}

\end{document}
